% Lecture 2: Concurrency: event loop, isolates, and goroutines
% Compile: latexmk -xelatex lecture02.tex
\documentclass[10pt,aspectratio=169,t]{beamer}
\usetheme{upbminimal}

\course{Mobile and Embedded Computing}
\decklabel{Lecture 2}
\title{Concurrency}
\subtitle{Event loop, isolates, and goroutines}
\author{Dinu-Ștefan Rusu}
\institute{FILS, Universitatea Politehnica București}
\date{Spring 2027}

\begin{document}

\titleframe

\objectivesframe{
  \cell[\faClock]{Respect the frame budget}{One frame is 1 \(\div\) refresh rate. Blocking the UI thread past that costs a dropped frame, not a crash.}
  \cell[\faIcon{random}]{Move CPU work off the UI thread}{compute() for a one-shot call, Isolate.spawn for a worker that lives longer or streams results back.}
  \cell[\faIcon{share-alt}]{Write goroutines and channels}{Go's concurrency primitives: go func(), typed channels, and select, on the server side of the same stack.}
  \cell[\faIcon{exchange-alt}]{Compare the two models}{Dart isolates copy messages between heaps. Goroutines share one heap and rely on channels for discipline.}
}

\divider{Part 1 \textperiodcentered\ Foundations}{The frame budget and the event loop}[One thread, one queue, and what running on it costs]

\begin{frame}[eyebrow=Foundations]{One frame budget \(=\) 1 \(\div\) refresh rate}
  \vskip 2mm
  \begin{grid}[3]
    \bignum{16.7 ms}{60 Hz: older devices, power-saving mode}
    \bignum{11.1 ms}{90 Hz: common mid-range in 2026}
    \bignum{8.3 ms}{120 Hz: most 2026 flagships}
  \end{grid}
  \vskip 4mm
  \begin{itemize}
    \item Inside that budget the UI thread has to run build, layout and paint, then hand the frame to the GPU.
    \item Miss the budget and the frame is dropped. Users see that as a single stutter; repeated overruns look like continuous jank.
    \item Do not hard-code 16 ms. Phones change refresh rate at runtime to save power, so the budget itself moves.
  \end{itemize}
  \note{The old assumption was a fixed 16 ms budget. Most 2026 phones run 90 to 120 Hz, so the real number is 1 divided by whatever refresh rate the phone is using right now. Ask the room what happens to the budget when a phone drops to 60 Hz to save battery: it gets easier to hit, not harder.}
\end{frame}

\begin{frame}[eyebrow=Foundations]{The event loop and the microtask queue}
  \begin{columns}[T]
    \begin{column}{0.56\textwidth}
      \begin{itemize}
        \item Your app runs on a single thread inside the main isolate, and that isolate owns the UI.
        \item That thread runs an event loop: run one event, drain the entire microtask queue, then take the next event.
        \item \code{await} starts nothing itself. It registers a continuation and hands control back to the loop.
        \item Nothing else runs meanwhile: a 400 ms function costs 400 ms of dropped frames.
      \end{itemize}
    \end{column}
    \begin{column}{0.40\textwidth}
      \begin{enumerate}
        \item Microtask queue: Future callbacks, drained completely first.
        \item Event queue: I/O, timers, gestures, platform messages.
        \item Frame: build, then layout, then paint.
      \end{enumerate}
    \end{column}
  \end{columns}
  \begin{takeaway}{Concurrent, not parallel.}
    One thread, interleaving work that is mostly waiting.
  \end{takeaway}
  \note{Draw the loop on the board if the room looks unconvinced: one event, drain every microtask, then the next event, in that strict order. This is also why a synchronous exception inside a callback can starve the frame: nothing else gets a turn until that callback returns.}
\end{frame}

\begin{frame}[fragile,eyebrow=Foundations]{Futures and async/await: a refresher}
  \begin{columns}[T]
    \begin{column}{0.56\textwidth}
\begin{dart}
Future<User> loadUser(String id) async {
  final res = await dio.get('/users/$id');
  return User.fromJson(res.data);
}

Future<void> refresh() async {
  try {
    user = await loadUser(id);
  } finally {
    if (mounted) setState(() {});
  }
}
\end{dart}
    \end{column}
    \begin{column}{0.40\textwidth}
      \lead{ADMD Lecture 4 owns this topic.}{This slide is a reminder, not new material.}
      \muted{A failure inside async code completes the Future with an error instead of throwing at the call site. \code{await} rethrows it, so try/catch behaves exactly as it does in synchronous code.}
    \end{column}
  \end{columns}
  \note{This is deliberately short. ADMD Lecture 4 already taught Future, async, await and error handling in depth. The only new point worth making here is the connection to the event loop: await is what hands control back to the loop covered on the previous slide.}
\end{frame}

\begin{frame}[fragile,eyebrow=Foundations]{Streams: many values, over time}
  \begin{columns}[T]
    \begin{column}{0.56\textwidth}
\begin{dart}
// A Stream is many values, over time.
Stream<int> ticker() async* {
  for (var i = 0; ; i++) {
    await Future.delayed(const Duration(seconds: 1));
    yield i;
  }
}
\end{dart}
    \end{column}
    \begin{column}{0.40\textwidth}
      \begin{itemize}
        \item A Future is one value, later. A Stream is many values, over time.
        \item StreamBuilder subscribes, rebuilds on each event and cancels for you.
        \item Cancel every subscription you create by hand, in \code{dispose()}.
      \end{itemize}
    \end{column}
  \end{columns}
  \begin{takeaway}{One value gives a Future. Many values give a Stream.}
    Both ride the same event loop. \(\rightarrow\) WebSockets and sensor feeds in later lectures.
  \end{takeaway}
  \note{Point out that async* and yield are to Stream what async and return are to Future. The consumer side, await for, is also just sugar over the same event loop: it still processes one event and drains microtasks before the next.}
\end{frame}

\begin{frame}[fragile,eyebrow=Foundations]{Async is not parallel}
  \begin{columns}[T]
    \begin{column}{0.56\textwidth}
\begin{dart}
// This function is async. It still
// freezes the UI.
Future<int> sumTo(int n) async {
  var total = 0;
  for (var i = 0; i < n; i++) {
    total += i;   // pure CPU, no yield
  }
  return total;   // the loop never turned
}
\end{dart}
    \end{column}
    \begin{column}{0.40\textwidth}
      \begin{itemize}
        \item \code{async} only helps when something else does the waiting: a socket, a disk, a timer.
        \item Pure CPU work never returns to the event loop. Cooperative multitasking only works if every task yields.
      \end{itemize}
      \muted{Typical offenders: a large JSON decode, image resize, encryption, sorting a big list.}
    \end{column}
  \end{columns}
  \begin{takeaway}{Rule of thumb.}
    Waiting on I/O: \code{await}. Burning CPU for longer than a frame: an isolate.
  \end{takeaway}
  \note{This is the slide that motivates the whole next section. Ask the room to spot why sumTo is misleading: it is marked async, but the keyword changes nothing about how the loop body executes. The fix is not a better await, it is moving the loop to an isolate.}
\end{frame}

\divider{Part 2 \textperiodcentered\ Isolates}{Moving CPU work off the UI thread}[A second heap, not a second thread inside yours]

\begin{frame}[fragile,eyebrow=Isolates]{compute(): the one-call isolate}
  \begin{columns}[T]
    \begin{column}{0.56\textwidth}
\begin{dart}
import 'package:flutter/foundation.dart';

// Must be a top-level or static function.
List<Photo> parsePhotos(String body) {
  final list = jsonDecode(body) as List;
  return list.map(Photo.fromJson).toList();
}

Future<List<Photo>> fetchPhotos() async {
  final res = await http.get(url);   // I/O
  return compute(parsePhotos, res.body);
}
\end{dart}
    \end{column}
    \begin{column}{0.40\textwidth}
      \begin{itemize}
        \item \code{compute()} spawns a short-lived isolate, runs one function, sends the result back and shuts it down.
        \item Arguments and results must be sendable: primitives, Strings, Lists and Maps of them.
        \item It costs a few milliseconds to spawn and copies the payload both ways. Do not reach for it to save 2 ms.
        \item Outside Flutter, \code{Isolate.run} does the same job in plain Dart.
      \end{itemize}
    \end{column}
  \end{columns}
  \note{compute() covers most isolate work application code actually needs. The next slide covers the cases it does not: a worker that stays alive, or one that streams several results back instead of returning once.}
\end{frame}

\begin{frame}[fragile,eyebrow=Isolates]{Isolate.spawn, SendPort, ReceivePort}
  \begin{columns}[T]
    \begin{column}{0.58\textwidth}
\begin{dart}
Future<int> sumInIsolate(int n) async {
  final inbox = ReceivePort();
  await Isolate.spawn(_worker, [inbox.sendPort, n]);
  return await inbox.first as int;
}

void _worker(List msg) {
  Isolate.exit(msg[0], sumTo(msg[1] as int));
}
\end{dart}
    \end{column}
    \begin{column}{0.38\textwidth}
      \begin{itemize}
        \item ReceivePort is your inbox. SendPort is a handle other isolates can post to.
        \item Messages are deep-copied between heaps: nothing is shared, so nothing needs a lock.
        \item \code{Isolate.exit} hands the result over without copying it, then ends the isolate. \code{sumTo} is the loop from the previous slide.
      \end{itemize}
    \end{column}
  \end{columns}
  \note{Use spawn over compute() when the worker is long-lived or streams several results back instead of one. Point out that the worker function only ever sees what arrived in msg: it has no access to any variable from the isolate that spawned it.}
\end{frame}

\begin{frame}[eyebrow=Isolates]{Async vs isolates, precisely}
  \vskip 1mm
  {\usebeamerfont{small}\begin{tabular}{@{}p{22mm}p{54mm}p{54mm}@{}}
    \toprule
    & \thead{async / await} & \thead{Isolate} \\
    \midrule
    Threads        & one, on the UI thread                 & one per isolate, scheduled by the OS \\
    Memory         & one heap, shared by everything         & a separate heap per isolate \\
    Talking        & just call the function                 & messages over ports, copied \\
    Good for       & I/O: network, disk, database, timers    & CPU: parse, decode, encrypt, sort \\
    Data races     & impossible: one thing runs at a time    & impossible: no shared memory \\
    Ordering races & possible                                & possible: replies can reorder \\
    \bottomrule
  \end{tabular}}
  \begin{takeaway}{No shared memory means no data races. It does not mean no races.}
    A reply can still arrive out of order, and isolate count is not capped by cores.
  \end{takeaway}
  \note{Two corrections to keep in mind. Race conditions are not impossible with isolates, only data races are: logical and message-ordering races remain possible. And spawning more isolates than you have cores is fine; you simply stop gaining throughput once the cores are saturated.}
\end{frame}

\divider{Part 3 \textperiodcentered\ Goroutines}{Go's concurrency model}[Shared memory, channels, and one goroutine per request]

\begin{frame}[fragile,eyebrow=Goroutines]{Why Go for the server side of this stack}
  \begin{columns}[T]
    \begin{column}{0.56\textwidth}
      \begin{itemize}
        \item \textbf{One static binary.} \code{go build} produces one executable, no runtime and no dependency folder to ship alongside it.
        \item \textbf{It starts in milliseconds.} Nothing to boot means short cold starts on a scale-to-zero platform.
        \item \textbf{The standard library is already the server.} \code{net/http} and \code{encoding/json} need no framework.
        \item The goroutines this section covers are one per request, by default, in that same standard library.
      \end{itemize}
    \end{column}
    \begin{column}{0.40\textwidth}
\begin{shell}
CGO_ENABLED=0 GOOS=linux \
  GOARCH=arm64 \
  go build -o server ./cmd/server

scp server root@vps:/usr/local/bin/
ssh root@vps 'systemctl restart api'
\end{shell}
    \end{column}
  \end{columns}
  \note{This is a preview, not the full backend picture: Lecture 10 covers serverless versus a VPS and deploying Go in depth. Here the point is narrower: Go's standard library gives you a working concurrent server before you write a single goroutine yourself.}
\end{frame}

\begin{frame}[eyebrow=Goroutines]{Same stack, a different runtime}
  \vskip 2mm
  \begin{grid}[2]
    \cell[\faIcon{mobile-alt}]{On the phone}{One thread owns the UI, and it has a frame budget to protect. Concurrency exists to keep that one thread free.}
    \cell[\faServer]{On the server}{No frame to protect, only requests to serve. Concurrency exists to handle many of them at once, on every core you have.}
  \end{grid}
  \vskip 5mm
  \begin{itemize}
    \item The course stack is Flutter and Dart on the phone, Go on the server. Same problem, concurrency, two different reasons to reach for it.
    \item Go's answer is the goroutine: a function that runs concurrently with the rest of the program, cheap enough to start one per incoming request.
  \end{itemize}
  \note{Frame this as the pivot of the lecture: everything up to now protected one deadline on one device. The rest of the lecture is about a runtime built to serve many independent callers at once, with no deadline of its own.}
\end{frame}

\begin{frame}[fragile,eyebrow=Goroutines]{A Go handler: one goroutine per request}
  \begin{columns}[T]
    \begin{column}{0.56\textwidth}
\begin{golang}
type Reading struct {
    DeviceID string  `json:"device_id"`
    TempC    float64 `json:"temp_c"`
}
func handleReading(w http.ResponseWriter,
    r *http.Request) {
    var in Reading
    json.NewDecoder(r.Body).Decode(&in)
    w.WriteHeader(http.StatusCreated)
}
\end{golang}
    \end{column}
    \begin{column}{0.40\textwidth}
      \begin{itemize}
        \item \code{net/http} hands every connection to its own goroutine automatically.
        \item \code{handleReading} already runs concurrently for every caller. Nothing here asked for that.
        \item \code{"POST /readings"} routing is built into the standard library, no framework needed.
      \end{itemize}
      \textbf{The goroutine is what this section explains next.}
    \end{column}
  \end{columns}
  \note{Live demo if time allows: go run the handler, then fire two curl requests back to back and log the goroutine id inside the handler to show they differ. The framework did nothing special; this is the default behavior of net/http.}
\end{frame}

\begin{frame}[fragile,eyebrow=Goroutines]{Goroutines: cheap, runtime-scheduled}
  \begin{columns}[T]
    \begin{column}{0.56\textwidth}
\begin{golang}
func main() {
    for i := 0; i < 5; i++ {
        go worker(i)
    }
    time.Sleep(time.Second)
}

func worker(id int) {
    fmt.Println("worker", id)
}
\end{golang}
    \end{column}
    \begin{column}{0.40\textwidth}
      \begin{itemize}
        \item \code{go} returns at once; it never waits for the call.
        \item A goroutine starts at 2 KB of stack. An OS thread starts at megabytes.
        \item The runtime spreads goroutines over a few OS threads, one per core.
      \end{itemize}
    \end{column}
  \end{columns}
  \begin{takeaway}{Cheap enough to have thousands.}
    An idle goroutine costs almost nothing.
  \end{takeaway}
  \note{The Sleep here is a placeholder, not a pattern: main exiting kills every goroutine still running, sleeping is just a way to give them a chance to print before that happens. The next slide replaces it with the correct tool.}
\end{frame}

\begin{frame}[fragile,eyebrow=Goroutines]{Waiting for many goroutines: sync.WaitGroup}
  \begin{columns}[T]
    \begin{column}{0.56\textwidth}
\begin{golang}
var wg sync.WaitGroup
for i := 0; i < 5; i++ {
    wg.Add(1)
    go func(id int) {
        defer wg.Done()
        process(id)
    }(i)
}
wg.Wait()
\end{golang}
    \end{column}
    \begin{column}{0.40\textwidth}
      \begin{itemize}
        \item \code{Add} before the goroutine starts, \code{Done} when it finishes, usually deferred.
        \item \code{Wait} blocks the caller until the count reaches zero.
        \item Recent Go gives each iteration its own \code{id}, so passing it is habit now, not a fix.
      \end{itemize}
    \end{column}
  \end{columns}
  \begin{takeaway}{This replaces the Sleep from the previous slide.}
    Wait for the actual work to finish, not for a guessed amount of time.
  \end{takeaway}
  \note{WaitGroup answers only one question: has everyone finished. It carries no data back. The next few slides cover channels, which is how goroutines actually exchange values while they run.}
\end{frame}

\begin{frame}[fragile,eyebrow=Goroutines]{Channels: typed, synchronized queues}
  \begin{columns}[T]
    \begin{column}{0.56\textwidth}
\begin{golang}
ch := make(chan int) // unbuffered
go func() {
    ch <- 42 // blocks until read
}()
value := <-ch // blocks until sent
close(ch) // no more sends
\end{golang}
    \end{column}
    \begin{column}{0.40\textwidth}
      \begin{itemize}
        \item A channel is a typed pipe between goroutines. Send and receive block by default, and that is what synchronizes them.
        \item Closing a channel signals ``no more values''. A receive from a closed channel returns immediately.
        \item Only the sender should close a channel.
      \end{itemize}
    \end{column}
  \end{columns}
  \begin{takeaway}{Rob Pike, Go proverb.}
    ``Don't communicate by sharing memory; share memory by communicating.''
  \end{takeaway}
  \note{Contrast this immediately with Dart: Go shares one heap and uses the channel as a convention to make that safe. Nothing in the language stops a goroutine from touching the variable directly instead; the discipline is on the programmer, not enforced like Dart's separate heaps.}
\end{frame}

\begin{frame}[fragile,eyebrow=Goroutines]{Buffered channels and backpressure}
  \begin{columns}[T]
    \begin{column}{0.56\textwidth}
\begin{golang}
ch := make(chan int, 3) // cap 3
ch <- 1
ch <- 2
ch <- 3 // buffer not full yet
ch <- 4 // buffer full: blocks
\end{golang}
    \end{column}
    \begin{column}{0.40\textwidth}
      \begin{itemize}
        \item A buffered channel decouples sender and receiver up to \(N\) values.
        \item Send blocks only once the buffer is full. Receive blocks only once it is empty.
        \item Picking \(N\) is a capacity decision, the same trade-off as sizing any queue.
      \end{itemize}
    \end{column}
  \end{columns}
  \begin{takeaway}{A bounded queue, for free.}
    No separate library: the channel itself gives you backpressure.
  \end{takeaway}
  \note{Ask what happens if N is set far too high: the program stops applying any backpressure at all, and a slow consumer allows memory to grow without bound. Buffering trades backpressure for a small amount of slack; it does not remove the need for one or the other.}
\end{frame}

\begin{frame}[eyebrow=Goroutines]{Goroutine and channel pitfalls}
  \vskip 1mm
  {\usebeamerfont{small}\begin{tabular}{@{}p{28mm}p{47mm}p{47mm}@{}}
    \toprule
    \thead{Pitfall} & \thead{What happens} & \thead{Fix} \\
    \midrule
    Deadlock & Every goroutine blocks on a channel with no partner to unblock it & Add a receiver, a buffer, or a \code{select} with a default \\
    Goroutine leak & A goroutine blocks forever on a channel nobody sends on or closes & Pass a context and select on its \code{Done()} channel \\
    Data race & Two goroutines read and write one variable with no channel or mutex & Run \code{go run -race}, or guard the variable with \code{sync.Mutex} \\
    \bottomrule
  \end{tabular}}
  \begin{takeaway}{None of this is checked for you.}
    The Dart isolate table two sections back had two rows of ``impossible''. Go has none: shared memory buys speed, and hands the safety work back to you.
  \end{takeaway}
  \note{This is the honest cost of Go's model. A data race here is a real, silent bug: two goroutines touching the same map with no synchronization can corrupt it without ever panicking. The race detector is not optional in a codebase with real concurrency; run it in CI.}
\end{frame}

\begin{frame}[fragile,eyebrow=Goroutines]{select: waiting on multiple channels}
  \begin{columns}[T]
    \begin{column}{0.56\textwidth}
\begin{golang}
select {
case v := <-ch1:
    fmt.Println("from ch1:", v)
case v := <-ch2:
    fmt.Println("from ch2:", v)
}
\end{golang}
    \end{column}
    \begin{column}{0.40\textwidth}
      \begin{itemize}
        \item \code{select} blocks until one of its cases can proceed.
        \item If several are ready at once, Go picks one at random, so no single channel is starved.
        \item Dart has no direct equivalent: the closest is \code{Future.any}, which races Futures.
      \end{itemize}
    \end{column}
  \end{columns}
  \note{This is the piece that makes Go's channels a language feature rather than just a queue type. Dart's Future.any comes closest, but select is built into the language and works uniformly across any number of channels.}
\end{frame}

\begin{frame}[fragile,eyebrow=Goroutines]{select with a timeout and a default case}
  \begin{columns}[T]
    \begin{column}{0.56\textwidth}
\begin{golang}
select {
case v := <-ch:
    fmt.Println(v)
case <-time.After(2 * time.Second):
    fmt.Println("timed out")
default:
    fmt.Println("nothing ready yet")
}
\end{golang}
    \end{column}
    \begin{column}{0.40\textwidth}
      \begin{itemize}
        \item \code{time.After} returns a channel that fires once, after the given duration: racing it against \code{ch} is a timeout.
        \item \code{default} makes the whole \code{select} non-blocking: nothing ready, run this branch and move on.
        \item Never combine a real case with \code{default} unless you mean ``check once and give up''.
      \end{itemize}
    \end{column}
  \end{columns}
  \begin{takeaway}{select turns waiting into a choice.}
    Exactly one branch runs; nothing else blocks the goroutine.
  \end{takeaway}
  \note{Point out that this timeout pattern is idiomatic Go: there is no separate Future.timeout call, the same select statement that multiplexes channels also expresses a deadline. It is one primitive doing both jobs.}
\end{frame}

\divider{Part 4 \textperiodcentered\ Compared}{Producer and consumer, Dart and Go}[Same shape, opposite guarantee]

\begin{frame}[fragile,eyebrow=Compared]{Producer/consumer in Dart}
\begin{dart}
final inbox = ReceivePort();
await Isolate.spawn(producer, inbox.sendPort);
await for (final item in inbox) {   // consumer
  if (item == null) break;
  print(item);
}
void producer(SendPort out) {
  for (var i = 0; i < 5; i++) out.send(i);
  out.send(null);
}
\end{dart}
  \begin{takeaway}{The two heaps never touch.}
    Every value crossing \code{inbox} is copied. \code{null} is this example's own end-of-stream marker; the loop reads it, the language does not.
  \end{takeaway}
  \note{The consumer treats inbox, a ReceivePort, as a Stream: await for works on it directly. The producer runs in a separate isolate and can only reach the consumer through the SendPort it was handed at spawn time.}
\end{frame}

\begin{frame}[fragile,eyebrow=Compared]{Producer/consumer in Go}
\begin{golang}
ch := make(chan int)
go func() {            // producer
    for i := 0; i < 5; i++ {
        ch <- i
    }
    close(ch)
}()
for item := range ch { // consumer
    fmt.Println(item)
}
\end{golang}
  \begin{takeaway}{The heap is shared.}
    Nothing stops another goroutine from reading \code{ch} too: a convention, not a guarantee.
  \end{takeaway}
  \note{range over a channel receives until it is closed, then stops on its own: no ad hoc null marker needed, unlike the Dart version on the facing slide. Both goroutines here run on the same heap; the channel is the only synchronization between them.}
\end{frame}

\begin{frame}[eyebrow=Compared]{Dart isolates vs Go goroutines}
  \vskip 1mm
  {\usebeamerfont{small}\begin{tabular}{@{}p{24mm}p{53mm}p{53mm}@{}}
    \toprule
    & \thead{Dart isolate} & \thead{Go goroutine} \\
    \midrule
    Memory        & a separate heap per isolate  & one heap, shared by every goroutine \\
    Talking       & ports; messages are copied   & channels; you choose what crosses \\
    Start-up cost & a few ms, a fresh heap       & about 2 KB of stack, growing \\
    Typical count & dozens, spawned deliberately & thousands, one per connection \\
    Data races    & impossible: no shared memory & possible: guard it yourself \\
    On failure    & ends that isolate only       & an unrecovered panic ends the process \\
    \bottomrule
  \end{tabular}}
  \begin{takeaway}{Same task, opposite guarantee.}
    Isolates trade a copy per message for safety. Goroutines trade the copy for speed, and hand the discipline back to you.
  \end{takeaway}
  \note{Go's slogan is a convention, not an enforced rule: nothing stops a goroutine from reaching across the shared heap instead of using the channel. Dart's separate heaps make that mistake impossible to write, at the cost of copying every message.}
\end{frame}

\begin{frame}[eyebrow=Compared]{What each model buys you: phone vs server}
  \vskip 2mm
  \begin{grid}[2]
    \cell[\faIcon{mobile-alt}]{On the phone}{Protect one frame budget. Anything over a few milliseconds of CPU work moves to an isolate; you spawn only as many as the work actually needs.}
    \cell[\faServer]{On the server}{Maximize throughput, with no frame to protect. One goroutine per request is the default because it is cheap, spread across every core with no extra code from you.}
  \end{grid}
  \begin{takeaway}{Concurrency solves a different problem on each side.}
    The phone spends it protecting a deadline. The server spends it on throughput. \(\rightarrow\) Lecture 10 puts a whole Go backend on top of this model.
  \end{takeaway}
  \note{Close the comparison by asking the room: if a mobile screen needed to serve a thousand simultaneous callers, would isolates be the right tool? No: that is exactly the server-side problem Go was built for, and it is why the course teaches both.}
\end{frame}

\begin{frame}[eyebrow=Recap]{Three things to remember}
  \vskip 2mm
  \begin{enumerate}
    \item One frame budget is 1 divided by refresh rate. \muted{The event loop runs one thing at a time and drains every microtask before the next event.}
    \item compute() and Isolate.spawn move CPU work to a separate heap. \muted{Messages are copied, so data races are impossible, but ordering races are not.}
    \item Goroutines share one heap and are synchronized with channels and select. \muted{Cheap enough to use one per request, but a race or a panic is now your job to prevent.}
  \end{enumerate}
  \vskip 5mm
  \kv{Further reading}{\link{https://dart.dev/language/concurrency}{dart.dev/language/concurrency} \textperiodcentered\ \link{https://go.dev/tour/concurrency}{go.dev/tour/concurrency} \textperiodcentered\ \link{https://go.dev/doc/effective_go}{go.dev/doc/effective\_go} \textperiodcentered\ \link{https://api.flutter.dev}{api.flutter.dev: Isolate}}
  \note{Leave the room with the one-line version: async is for waiting, isolates are for a phone's CPU work, goroutines are for a server's request volume. Lecture 10 builds the Go side of the course project on exactly the goroutine-per-request pattern from this lecture.}
\end{frame}

\closingframe{Questions?}{dinu\_stefan.rusu@upb.ro \textperiodcentered\ Microsoft Teams: course channel}

\end{document}
